\documentclass[10pt,a4paper]{amsart}
\usepackage{amsmath,amssymb,amsthm,graphicx}
% Retained preamble for the Stacks Project chapter snapshot in this directory.
% The chapter's own preamble.tex is a copy of the shared Stacks Project preamble; the block
% below is sliced MECHANICALLY, line for line, from it. Nothing was added, removed or
% redefined; the wrapper class, the xr-hyper/hyperref wiring and the externaldocument list
% are omitted because the assembler supplies the document class.
%
% lines taken from preamble.tex: 3, 5, 9, 11, 12, 156, 157, 158, ... (51 lines total)
%
% package coverage restored from the chapter's own preamble (2026-10-13): xy, xy-2cell

\usepackage[all]{xy}
\xyoption{2cell}
\UseAllTwocells

\newenvironment{reference}{\comment}{\endcomment}
\newenvironment{slogan}{\comment}{\endcomment}
\newenvironment{history}{\comment}{\endcomment}
\theoremstyle{plain}
\newtheorem{theorem}[subsection]{Theorem}
\newtheorem{proposition}[subsection]{Proposition}
\newtheorem{lemma}[subsection]{Lemma}
\theoremstyle{definition}
\newtheorem{definition}[subsection]{Definition}
\newtheorem{example}[subsection]{Example}
\newtheorem{exercise}[subsection]{Exercise}
\newtheorem{situation}[subsection]{Situation}
\theoremstyle{remark}
\newtheorem{remark}[subsection]{Remark}
\newtheorem{remarks}[subsection]{Remarks}
\numberwithin{equation}{subsection}
\def\lim{\mathop{\mathrm{lim}}\nolimits}
\def\colim{\mathop{\mathrm{colim}}\nolimits}
\def\Spec{\mathop{\mathrm{Spec}}}
\def\Hom{\mathop{\mathrm{Hom}}\nolimits}
\def\Ext{\mathop{\mathrm{Ext}}\nolimits}
\def\SheafHom{\mathop{\mathcal{H}\!\mathit{om}}\nolimits}
\def\SheafExt{\mathop{\mathcal{E}\!\mathit{xt}}\nolimits}
\def\Sch{\mathit{Sch}}
\def\Mor{\mathop{\mathrm{Mor}}\nolimits}
\def\Ob{\mathop{\mathrm{Ob}}\nolimits}
\def\Sh{\mathop{\mathit{Sh}}\nolimits}
\def\NL{\mathop{N\!L}\nolimits}
\def\CH{\mathop{\mathrm{CH}}\nolimits}
\def\proetale{{pro\text{-}\acute{e}tale}}
\def\etale{{\acute{e}tale}}
\def\QCoh{\mathit{QCoh}}
\def\Ker{\mathop{\mathrm{Ker}}}
\def\Im{\mathop{\mathrm{Im}}}
\def\Coker{\mathop{\mathrm{Coker}}}
\def\Coim{\mathop{\mathrm{Coim}}}
\DeclareMathSymbol{\boxtimes}{\mathbin}{AMSa}{"02}
\def\QCohstack{\mathcal{QC}\!\mathit{oh}}
\def\Cohstack{\mathcal{C}\!\mathit{oh}}
\def\Spacesstack{\mathcal{S}\!\mathit{paces}}
\def\Quotfunctor{\mathrm{Quot}}
\def\Hilbfunctor{\mathrm{Hilb}}
\def\Curvesstack{\mathcal{C}\!\mathit{urves}}
\def\Polarizedstack{\mathcal{P}\!\mathit{olarized}}
\def\Complexesstack{\mathcal{C}\!\mathit{omplexes}}
\def\Pic{\mathop{\mathrm{Pic}}\nolimits}
\def\Picardstack{\mathcal{P}\!\mathit{ic}}
\def\Picardfunctor{\mathrm{Pic}}
\def\Deformationcategory{\mathcal{D}\!\mathit{ef}}
\usepackage{cleveref}
\crefname{originallemma}{original lemma}{original lemmas}
\Crefname{originallemma}{Original lemma}{Original lemmas}
\makeatletter
\newlabel{example-graded-algebras}{{0D3L}{source}{}{sourcepointer-1}{}}
\newlabel{example-graded-algebras@cref}{{[originallemma][0][]0D3L}{[1][0][]source}{}{}{}}
\makeatother
\title{Stacks, Examples of Deformations: lemma-graded-algebras-TI}
\author{The Stacks Project Authors}
\date{stacks-project:examples-defos.tex:local-snapshot-20261012; GNU FDL 1.2 or later}
\begin{document}
\maketitle
\noindent\textit{Editorial note.} Statement span is lines 880-890 of sources/stacks-examples-defos/pinned-original-source.tex; proof span is lines 892-984, the complete author proof (91 body lines).\par\medskip
\noindent\textit{Editorial note.} Navigation stubs: 1 cited-but-not-included labels are supplied as non-load-bearing pointers with their pinned tags where the snapshot tag index has them.\par\medskip
\noindent\textit{Editorial note.} extra\_packages is empty on purpose: the retained preamble already loads xy with options and a second declaration clashes in the assembler wrapper.\par\medskip
\noindent\textit{Editorial note.} Text is retained verbatim, including the author's own typos and placeholders.\par\medskip
\noindent\textit{Editorial note.} Queue index 2 of queue\_examples-defos.json (pb=3188, stmt 880-890, proof 892-984).\par\medskip
\noindent\textit{Allowed standard prerequisites (editorial).}\begin{itemize}
\item Explicit prerequisite boundary. The cited inputs are prerequisites; the argument in the retained span is not replaced by any cited result.
\end{itemize}\medskip
\ifdefined\hypertarget\hypertarget{sourcepointer-1}{}\fi
\noindent\textit{Source reference.} Original lemma 0D3L: sources/stacks-examples-defos/pinned-original-source.tex, source lines 825--835. Cited by the retained statement or proof but not part of the collected spans; supplied only as a navigation pointer. The argument does not depend on its content.\par\medskip
\begin{lemma}
\label{lemma-graded-algebras-TI}
In Example \ref{example-graded-algebras} let $P$ be a graded $k$-algebra.
Then
$$
T\Deformationcategory_P
\quad\text{and}\quad
\text{Inf}(\Deformationcategory_P) = \text{Der}_k(P, P)
$$
are finite dimensional if $P$ is finitely generated over $k$.
\end{lemma}

\begin{proof}
We first deal with the infinitesimal automorphisms.
Let $Q = P \otimes_k k[\epsilon]$.
Then an element of $\text{Inf}(\Deformationcategory_P)$
is given by an automorphism
$\gamma = \text{id} + \epsilon \delta : Q \to Q$
as above where now $\delta : P \to P$.
The fact that $\gamma$ is graded implies that
$\delta$ is homogeneous of degree $0$.
The fact that $\gamma$ is $k$-linear implies that
$\delta$ is $k$-linear.
The fact that $\gamma$ is multiplicative implies that
$\delta$ is a $k$-derivation.
Conversely, given a $k$-derivation $\delta : P \to P$
homogeneous of degree $0$, we obtain an automorphism
$\gamma = \text{id} + \epsilon \delta$ as above.
Thus we see that
$$
\text{Inf}(\Deformationcategory_P) = \text{Der}_k(P, P)
$$
as predicted in the lemma.
Clearly, if $P$ is generated in degrees $P_i$,
$0 \leq i \leq N$, then $\delta$ is determined by
the linear maps $\delta_i : P_i \to P_i$ for
$0 \leq i \leq N$ and we see that
$$
\dim_k \text{Der}_k(P, P) < \infty
$$
as desired.

\medskip\noindent
To finish the proof of the lemma we show that there is a finite
dimensional deformation space. To do this we
choose a presentation
$$
k[X_1, \ldots, X_n]/(F_1, \ldots, F_m) \longrightarrow P
$$
of graded $k$-algebras where $\deg(X_i) = d_i$ and
$F_j$ is homogeneous of degree $e_j$.
Let $Q$ be any graded $k[\epsilon]$-algebra
finite free in each degree which comes with an isomorphsm
$\alpha : Q/\epsilon Q \to P$ so that $(Q, \alpha)$ defines
an element of $T\Deformationcategory_P$.
Choose a homogeneous element $q_i \in Q$ of degree $d_i$
mapping to the image of $X_i$ in $P$.
Then we obtain
$$
k[\epsilon][X_1, \ldots, X_n] \longrightarrow Q,\quad
X_i \longmapsto q_i
$$
and since $P = Q/\epsilon Q$ this map is surjective by Nakayama's lemma.
A small diagram chase shows we can choose homogeneous elements
$F_{\epsilon, j} \in k[\epsilon][X_1, \ldots, X_n]$ of degree $e_j$
mapping to zero in $Q$ and mapping to $F_j$ in $k[X_1, \ldots, X_n]$.
Then
$$
k[\epsilon][X_1, \ldots, X_n]/(F_{\epsilon, 1}, \ldots, F_{\epsilon, m})
\longrightarrow Q
$$
is a presentation of $Q$ by flatness of $Q$ over $k[\epsilon]$.
Write
$$
F_{\epsilon, j} =  F_j + \epsilon G_j
$$
There is some ambiguity in the vector $(G_1, \ldots, G_m)$.
First, using different choices of $F_{\epsilon, j}$
we can modify $G_j$ by an arbitrary element of degree $e_j$
in the kernel of $k[X_1, \ldots, X_n] \to P$.
Hence, instead of $(G_1, \ldots, G_m)$, we remember the
element
$$
(g_1, \ldots, g_m) \in P_{e_1} \oplus \ldots \oplus P_{e_m}
$$
where $g_j$ is the image of $G_j$ in $P_{e_j}$.
Moreover, if we change our choice of $q_i$ into $q_i + \epsilon p_i$
with $p_i$ of degree $d_i$ then a computation (omitted) shows
that $g_j$ changes into
$$
g_j^{new} = g_j - \sum\nolimits_{i = 1}^n p_i \partial F_j / \partial X_i
$$
We conclude that the isomorphism class of $Q$ is determined by the
image of the vector $(G_1, \ldots, G_m)$ in the $k$-vector space
$$
W  = \Coker(P_{d_1} \oplus \ldots \oplus P_{d_n}
\xrightarrow{(\frac{\partial F_j}{\partial X_i})}
P_{e_1} \oplus \ldots \oplus P_{e_m})
$$
In this way we see that we obtain an injection
$$
T\Deformationcategory_P \longrightarrow W
$$
Since $W$ visibly has finite dimension, we conclude that the lemma is true.
\end{proof}

\end{document}
